% GENERATED FILE. Edit canonical lessons, then run npm run generate:books.
% canonical-source-hash: fnv1a64:ab8fc236dac60934
% canonical-lessons: LA-C222-aurum, LA-C222-argentum, LA-C222-aes, LA-C222-ferrum, LA-C222-plumbum

\chapter{Gold, Silver, Bronze, Iron, Metal lead}
\label{ch:la-gold-silver-bronze-iron-metal-lead}
\hlchaptermodality{222}

\begin{chapteropening}
\textbf{By the end of this chapter:} \emph{I can say gold, silver, bronze, iron and the metal lead.}

Complete the last lesson of chapter 222: I can say gold, silver, bronze, iron and the metal lead.
\end{chapteropening}

\section[aurum, aurī]{aurum, aurī — gold}

\label{lesson:LA-C222-aurum}



\begin{quote}
Before the new one: say the Latin for fortune, then the Latin for an accident.
\end{quote}

\subsection*{You'll want to know: aurum, aurī}
\textbf{aurum, aurī} — \textquotedblleft{}gold\textquotedblright{}.

\begin{cousinweb}
\textbf{aurum} is inherited, usually traced to Proto-Indo-European \emph{*h2eh2usom}, \textquotedblleft{}a glow\textquotedblright{}, from \emph{*h2ews-}, \textquotedblleft{}to dawn, to grow light\textquotedblright{} (EtymDB 2.1), the root of \emph{aurōra}, \textquotedblleft{}dawn\textquotedblright{}. Festus records a Sabine \emph{ausum}, \textquotedblleft{}gold\textquotedblright{}, which keeps the older s (Lewis \& Short).
\end{cousinweb}

\begin{grammarlens}[title={What you've built}]
One more word for this chapter.
\end{grammarlens}

\subsection*{Your turn}
\begin{itemize}
\raggedright
  \item \emph{Say it:} \emph{aurum, aurī}
  \item \emph{Say it:} \emph{aurum, aurī}, once more
  \item \emph{Say it:} say \emph{aurum, aurī}
  \item \emph{Recall:} say the Latin for fortune, then the Latin for an accident, then say \emph{aurum, aurī} again
\end{itemize}

\begin{tcolorbox}[breakable,colback=teal!4,colframe=teal!35!black,title={Before you move on}]
What does aurum, aurī mean? (\textquotedblleft{}Gold\textquotedblright{}.) Say it once more.
\end{tcolorbox}
\section[argentum, argentī]{argentum, argentī — silver}

\label{lesson:LA-C222-argentum}



\begin{quote}
Before the new one: say the Latin for an accident, then the Latin for gold.
\end{quote}

\subsection*{You'll want to know: argentum, argentī}
\textbf{argentum, argentī} — \textquotedblleft{}silver\textquotedblright{}.

\begin{cousinweb}
\textbf{argentum} is inherited, usually traced to Proto-Indo-European \emph{*h2rgntom}, \textquotedblleft{}silver\textquotedblright{}, from \emph{*h2erg-}, \textquotedblleft{}white, shining\textquotedblright{} (EtymDB 2.1); Greek \emph{argēs}, \textquotedblleft{}white, bright\textquotedblright{}, is a cousin (Lewis \& Short).
\end{cousinweb}

\begin{grammarlens}[title={What you've built}]
One more word for this chapter.
\end{grammarlens}

\subsection*{Your turn}
\begin{itemize}
\raggedright
  \item \emph{Say it:} \emph{argentum, argentī}
  \item \emph{Say it:} \emph{argentum, argentī}, once more
  \item \emph{Say it:} say \emph{argentum, argentī}
  \item \emph{Recall:} say the Latin for an accident, then the Latin for gold, then say \emph{argentum, argentī} again
\end{itemize}

\begin{tcolorbox}[breakable,colback=teal!4,colframe=teal!35!black,title={Before you move on}]
What does argentum, argentī mean? (\textquotedblleft{}Silver\textquotedblright{}.) Say it once more.
\end{tcolorbox}
\section[aes, aeris]{aes, aeris — bronze; money}

\label{lesson:LA-C222-aes}



\begin{quote}
Before the new one: say the Latin for gold, then the Latin for silver.
\end{quote}

\subsection*{You'll want to know: aes, aeris}
\textbf{aes, aeris} — \textquotedblleft{}bronze; money\textquotedblright{}.

\begin{cousinweb}
\textbf{aes} is usually traced to Proto-Italic \emph{*aos}, \textquotedblleft{}copper, bronze\textquotedblright{} (EtymDB 2.1); Gothic \emph{aiz}, \textquotedblleft{}bronze\textquotedblright{}, and English \emph{ore} are cousins (Lewis \& Short).
\end{cousinweb}

\begin{grammarlens}[title={What you've built}]
One more word for this chapter.
\end{grammarlens}

\subsection*{Your turn}
\begin{itemize}
\raggedright
  \item \emph{Say it:} \emph{aes, aeris}
  \item \emph{Say it:} \emph{aes, aeris}, once more
  \item \emph{Say it:} say \emph{aes, aeris}
  \item \emph{Recall:} say the Latin for gold, then the Latin for silver, then say \emph{aes, aeris} again
\end{itemize}

\begin{tcolorbox}[breakable,colback=teal!4,colframe=teal!35!black,title={Before you move on}]
What does aes, aeris mean? (\textquotedblleft{}Bronze; money\textquotedblright{}.) Say it once more.
\end{tcolorbox}
\section[ferrum, ferrī]{ferrum, ferrī — iron}

\label{lesson:LA-C222-ferrum}



\begin{quote}
Before the new one: say the Latin for silver, then the Latin for bronze; money.
\end{quote}

\subsection*{You'll want to know: ferrum, ferrī}
\textbf{ferrum, ferrī} — \textquotedblleft{}iron\textquotedblright{}.

\begin{cousinweb}
\textbf{ferrum} has no sure origin; it is often thought to be a loan from a language of the eastern Mediterranean, perhaps through Etruscan, but the source is unknown.
\end{cousinweb}

\begin{grammarlens}[title={What you've built}]
One more word for this chapter.
\end{grammarlens}

\subsection*{Your turn}
\begin{itemize}
\raggedright
  \item \emph{Say it:} \emph{ferrum, ferrī}
  \item \emph{Say it:} \emph{ferrum, ferrī}, once more
  \item \emph{Say it:} say \emph{ferrum, ferrī}
  \item \emph{Recall:} say the Latin for silver, then the Latin for bronze; money, then say \emph{ferrum, ferrī} again
\end{itemize}

\begin{tcolorbox}[breakable,colback=teal!4,colframe=teal!35!black,title={Before you move on}]
What does ferrum, ferrī mean? (\textquotedblleft{}Iron\textquotedblright{}.) Say it once more.
\end{tcolorbox}
\section[plumbum, plumbī]{plumbum, plumbī — the metal lead}

\label{lesson:LA-C222-plumbum}



\begin{quote}
Before the new one: say the Latin for bronze; money, then the Latin for iron.
\end{quote}

\subsection*{You'll want to know: plumbum, plumbī}
\textbf{plumbum, plumbī} — \textquotedblleft{}the metal lead\textquotedblright{}.

\begin{cousinweb}
\textbf{plumbum} is related to Greek \emph{molybdos}, \textquotedblleft{}lead\textquotedblright{} (Lewis \& Short), but neither came straight from the other; both are thought to be loans from an unknown language. English \emph{plumb} and \emph{plumber} come from it.
\end{cousinweb}

\begin{grammarlens}[title={What you've built}]
One more word for this chapter.
\end{grammarlens}

\subsection*{Your turn}
\begin{itemize}
\raggedright
  \item \emph{Say it:} \emph{plumbum, plumbī}
  \item \emph{Say it:} \emph{plumbum, plumbī}, once more
  \item \emph{Say it:} say \emph{plumbum, plumbī}
  \item \emph{Recall:} say the Latin for bronze; money, then the Latin for iron, then say \emph{plumbum, plumbī} again
\end{itemize}

\begin{tcolorbox}[breakable,colback=teal!4,colframe=teal!35!black,title={Before you move on}]
What does plumbum, plumbī mean? (\textquotedblleft{}The metal lead\textquotedblright{}.) Say it once more.
\end{tcolorbox}
